\chapter{El otro circuito}

\label{cap:tierrafiscal}
\section{El segundo circuito}

Los capítulos anteriores siguieron sobre todo el mecanismo privado de producción de suelo. Es el dominante, pero no el único, y el capítulo \ref{cap:planteo} ya anunció el otro. Existe un segundo circuito, estatal y provincial, que distribuye tierra fiscal con finalidad social. Su órgano es la Secretaría de Tierras y Bienes del Estado, articulada con el Instituto Provincial de Vivienda, y su programa contemporáneo es el Plan Provincial ``Mi Lote''.

Un libro que describe la producción de suelo en un departamento periurbano y sólo mira al mercado privado deja fuera el instrumento con el que el Estado provincial interviene sobre el mismo objeto. Este capítulo lo documenta, y con él la primera tesis del libro queda completa.

\section{La cadena normativa}

\begin{ficha}{Régimen de tierra fiscal de la Provincia de Salta}
\begin{itemize}[leftmargin=*,nosep]
\item \textbf{Constitución provincial, art. 37:} los poderes públicos facilitan el acceso de los sectores de menores ingresos a una vivienda digna.
\item \textbf{Ley Nº 2616} (original 1338; sancionada el 24/08/1951, promulgada el 27/08/1951 y publicada en el Boletín Oficial Nº 4.025 del 03/09/1951; vigente: el Digesto Jurídico de la Ley 7913 la registra en su Anexo I, entre las leyes generales permanentes vigentes, fs.\ 90): autoriza a enajenar terrenos fiscales, \textbf{por adjudicación directa} los destinados a vivienda familiar, pequeños negocios o pequeñas quintas, y por concurso o licitación los destinados a comercio o industria. Su art. 2º fija el orden de preferencia: \textbf{1º los actuales ocupantes de los terrenos}; 2º los argentinos nativos; 3º quienes no tengan otro inmueble en la Provincia; 4º las personas de familia numerosa y limitados recursos; 5º por orden de presentación. El art. 5º, sustituido por la Ley 3590/1961, manda fijar el precio con un coeficiente de disminución del 30\%. El art. 6º prohíbe adjudicar más de dos parcelas por persona o familia.
\item \textbf{Decreto Nº 4.414/08} (B.O.\ Nº 17.974, 21/10/2008): declara la emergencia habitacional en todo el territorio provincial y crea el \textbf{Programa «Comunidad Organizada»}, en la Subsecretaría de Tierra y Hábitat, para «planificar, ordenar y transferir las tierras del dominio privado de la Provincia» a quienes «necesiten urgente solución habitacional», según la reglamentación que dictaría el Ministerio. Lo coordinan representantes de cinco organismos provinciales y «delegados de los Municipios involucrados en cada caso». Su artículo 4º \textbf{excluye del programa a quien esté incurso en infracciones vinculadas con la ocupación ilegal de inmuebles}.
\item \textbf{Decreto Nº 450/18}: encomienda al Instituto Provincial de la Vivienda la inscripción, asignación y preadjudicación de inmuebles del dominio privado de la Provincia a personas de menores ingresos.
\item \textbf{Ley Nº 7.905 y Decretos 18/15 y 21/15:} crean el Ministerio de Infraestructura, Tierra y Vivienda, la Secretaría de Tierra y Bienes y la Subsecretaría de Tierra y Hábitat.
\item \textbf{Decretos Nº 16/19 y 127/20:} con el cambio de gobierno, la Secretaría se vuelve a crear, ahora como «de Tierras y Bienes» y \textbf{bajo dependencia de la Secretaría General de la Gobernación}. El Decreto 16/19 (10/12/2019; B.O.\ Nº 20.649 del 18/12/2019; Sáenz--Posadas) aprueba las estructuras ministeriales, y su anexo ---organigramas publicados como imagen--- ubica la Secretaría de Tierras y Bienes del Estado, con una Subsecretaría de Regularización Dominial, dentro de la Secretaría General. El Decreto 127/20 (22/01/2020; B.O.\ Nº 20.672 del 27/01/2020), con vigencia desde el 11 de diciembre de 2019, aprueba la estructura de la Secretaría General hasta el nivel de dirección; su anexo, una copia certificada escaneada, cuelga de la Secretaría de Tierra y Bienes del Estado cuatro direcciones generales: \textbf{Inmuebles}, Tierras Fiscales y Bienes del Estado, Organización y Planificación, y la Unidad de Regulación Dominial ---así, «Regulación», en el anexo escaneado; el resto de los actos la llama de \textit{Regularización} Dominial, y es probable que la diferencia sea un error de transcripción del propio anexo---. No hay contradicción con las normas de 2015: son dos creaciones sucesivas del mismo organismo.
\item \textbf{Decreto Nº 399/20} (24/06/2020): crea el Plan Provincial de Urbanización Social y Regularización Dominial ``Mi Lote''.
\item \textbf{Resolución Nº 82/16} del Ministerio y su modificatoria \textbf{86/16}: Reglamento de Asignación de Lotes de Interés Social.
\item \textbf{Resoluciones Nº 7/16, 8/16 y 9/16} de la Secretaría: requisitos y documentación; tipologías de puntaje; procedimientos de asignación.
\item \textbf{Ley Nº 8274 y Decisión Administrativa Nº 19/2023} ---la de diciembre de 2023, dictada con la numeración del nuevo mandato; no debe confundirse con la Decisión Administrativa 19/23 de enero de ese año, que aprueba un contrato de servicios del Ministerio de Desarrollo Social---: la ley, modificatoria de la Ley de Ministerios 8171, asigna al Ministerio de Infraestructura la competencia de «entender en el desarrollo de las tierras fiscales» (art.\ 27, inc.\ 17); la decisión administrativa ---dictada el 26 de diciembre de 2023 por el Coordinador Administrativo de la Gobernación con refrendo del ministro Camacho, con vigencia desde el 11 de diciembre y publicada en el B.O.\ Nº 21.621 del 29/12/2023--- aprueba la estructura del ministerio y designa a sus autoridades. Su anexo, una copia certificada escaneada, ubica bajo el ministro tres secretarías ---Planificación, \textbf{Tierra y Bienes del Estado} y Obras Públicas---, con la \textbf{Subsecretaría de Regularización Dominial} dentro de la segunda, y la Dirección General de Inmuebles en la cabecera, dependiente directamente del ministro, junto al Instituto Provincial de Vivienda, Vialidad y la empresa de agua como entes vinculados. El texto de la Ley 8274 se conoce por los fundamentos del Decreto 198/24 (03/04/2024; B.O.\ Nº 21.684 del 08/04/2024) \pendiente{texto de la Ley 8274}.
\item \textbf{Decreto Nº 844/21} (27/09/2021): traslada la Secretaría al Ministerio de Infraestructura y dispone que la Dirección General de Inmuebles dependa directamente del Ministro.
\item \textbf{Dependencia actual:} la ficha oficial del organismo en el portal del Gobierno de Salta (consultada en septiembre de 2026) la ubica en la \textbf{Jefatura de Gabinete de Ministros}. El traslado lo dispuso la \textbf{Ley 8.511} (Decreto 789/25; B.O.\ Nº 22.080 del 26/11/2025), sancionada el 20 de noviembre de 2025 y vigente desde el 26, que derogó las Leyes 8171 y 8274 y \textbf{suprimió el Ministerio de Infraestructura}. Sus competencias pasaron a la Jefatura de Gabinete, que ahora «entiende» en el catastro y el registro inmobiliario, en el patrimonio del Estado y \textbf{en el desarrollo de las tierras fiscales} (art.\ 19, incs.\ 22 a 24), además de las obras públicas, la vivienda y las relaciones con los municipios. La estructura se aprobó por la Decisión Administrativa 569/25 (23/12/2025, vigente desde el 15 de diciembre; B.O.\ Nº 22.100): la Secretaría de Tierra y Bienes del Estado conserva su Subsecretaría de Regularización Dominial y sus direcciones generales de Tierras Fiscales y Bienes del Estado, de Organización y Planificación, de Regularización Dominial y de Gestión Financiera; la Dirección General de Inmuebles queda en la cabecera, bajo el jefe de Gabinete. Tampoco en esta estructura figura una unidad propia para el control de loteos. Hasta entonces, la Secretaría estaba en Infraestructura: el 3 de septiembre, la Decisión Administrativa 405/25 (B.O.\ Nº 22.027) todavía prorroga contratos de la Secretaría «en el ámbito del Ministerio de Infraestructura» hasta el 31 de diciembre de 2025, igual que antes las Decisiones 29/25 y 220/25 y los Decretos 64/23 y 198/24. En la segunda, el Decreto 503/26 (B.O.\ Nº 22.251), que la autoriza a firmar un comodato, ya lo dicta la Jefatura de Gabinete. Los actos de las dos etapas los refrenda el mismo funcionario, el \textbf{Ing.\ Sergio Darío Camacho}, cuya firma y sello de ministro de Infraestructura constan en el anexo de la Decisión Administrativa 29/25: al menos entre 2021 y 2025, como ministro de Infraestructura; en 2026, como jefe de Gabinete. El \textbf{Decreto 875/25}, publicado el 16 de diciembre de 2025, lo designa Jefe de Gabinete de Ministros. Y la Decisión Administrativa 80/26, del 11 de febrero de 2026 (B.O.\ Nº 22.133), aprueba desde el 1º de enero los contratos de la Secretaría «dependiente de la Jefatura de Gabinete de Ministros», con cargo a esa jurisdicción, invocando la Ley 8.511. \textbf{La Secretaría pasó con su ministro en el reordenamiento de fines de 2025}, el mismo en el que el organismo hídrico bajó a Subsecretaría y quedó también en la Jefatura de Gabinete (Decreto 52/26; capítulo \ref{cap:ribera}).
\end{itemize}
\end{ficha}

\begin{cautela}
\textbf{Sobre el nombre del organismo.} Los propios actos oficiales no lo escriben igual, y a veces difieren dentro del mismo año. Las normas de 2015 crean la «Secretaría de Tierra y Bienes». El Decreto 16/19 la vuelve a crear como «de Tierras y Bienes del Estado», el anexo del Decreto 127/20 la llama «de Tierra y Bienes del Estado» y la Resolución 466 D/20, «de Tierras y Bienes». En 2023, el Decreto 64/23 designa a su titular como «Secretaria de Tierras y Bienes del Estado» y el anexo de la Decisión Administrativa 19/23 dice «Tierra y Bienes del Estado»; la Decisión Administrativa 569/25 repite esta última forma, y la ficha oficial del portal provincial, consultada en septiembre de 2026, usa la primera, bajo la Jefatura de Gabinete de Ministros. \textbf{Este libro usa, en su propia prosa, «Secretaría de Tierra y Bienes» para 2015--2019 y «Secretaría de Tierras y Bienes del Estado» desde diciembre de 2019, que es la forma del portal vigente; en las citas y en las fichas conserva la grafía de cada instrumento.} La alternancia no indica organismos distintos: es el mismo, recreado y trasladado.
\end{cautela}

El artículo 2º del Decreto 844/21 merece atención por sí solo. Colocar a la Dirección General de Inmuebles bajo dependencia directa del Ministro de Infraestructura reúne, en una sola línea de mando, el organismo que confecciona y aprueba los planos de mensura y loteo --- el mismo que aprobó el plano 1.164 de El Durazno --- y el organismo que administra la tierra fiscal.

Los dos circuitos que este libro describe como paralelos convergen, en el organigrama, en el mismo despacho.

Y la convergencia es anterior a ese decreto. Entre diciembre de 2019 y septiembre de 2021, según el anexo del Decreto 127/20, la Dirección General de Inmuebles no estaba junto a la Secretaría de Tierras y Bienes del Estado: estaba dentro de ella, como una de sus cuatro direcciones generales. En ese mismo lapso, desde octubre de 2020, la Secretaría asumió además las competencias de la Unidad Coordinadora de Loteos (capítulo \ref{cap:loteo}). Durante casi dos años, una sola secretaría reunió la tierra fiscal, el registro y la aprobación de planos, y el control de los loteos privados. Es el período en que tramitó el expediente del plano de El Durazno, iniciado en 2019; la resolución que lo aprobó, en noviembre de 2021, la firmó ya el ministro de Infraestructura, dos meses después del traslado. El esquema de 2021 se mantuvo: la estructura aprobada en diciembre de 2023 (Decisión Administrativa 19/23) conserva a Inmuebles bajo el ministro y a la Secretaría de Tierras y Bienes del Estado a su lado, en el mismo ministerio. Y en ese organigrama la Unidad Coordinadora de Loteos no figura como unidad propia: sus competencias, encomendadas en 2020, no tienen un casillero visible.

La ficha oficial vigente lo dice sin rodeos. Entre las competencias publicadas de la Secretaría figuran «intervenir en todo lo relativo a la elaboración de planos y su posterior aprobación», «avalar las operaciones que realice la Dirección General de Inmuebles» y «analizar la pertinencia de requerir autorización legislativa para la disposición de bienes inmuebles». Lo que en 2021 fue una decisión de organigrama figura hoy como competencia propia del organismo que adjudica la tierra fiscal. La misma ficha, en el currículum de su titular, consigna su mandato como senadora provincial por el departamento La Caldera entre 2013 y 2021.

El decreto que encabeza la cadena merece dos observaciones. La primera es de tiempo: en octubre de 2008 declaró una emergencia y ordenó al Ministerio dictar la reglamentación del programa; el Reglamento de Asignación de Lotes de Interés Social es la Resolución 82/16, de septiembre de 2016. \textbf{Ocho años entre la emergencia y su reglamento}, y el reglamento dejó el contenido operativo en un anexo que el Boletín publica sólo como imagen escaneada, vinculada al acto.

La segunda es de origen. La cláusula que excluye a quienes estén incursos en infracciones por ocupación ilegal, y cuya tensión con la regularización de ocupantes este capítulo analiza a propósito del Plan Mi Lote, no nace en 2020: está en el decreto de 2008, fundada en que las ocupaciones «obstaculizan la adecuada planificación». Doce años después, el Decreto 399/20 la repite.

\section{Quién decide}

\begin{itemize}[leftmargin=*,nosep]
\item \textbf{Decreto 120/15} (21/12/2015): designa al C.P.N. Fernando Gabriel Martinis Mercado, Secretario de Tierra y Bienes. Es el firmante de las Resoluciones 7/16, 8/16 y 9/16.
\item \textbf{Decisión Administrativa 50/21} (29/01/2021): cesa a Pedro Guillermo Peretti por jubilación y designa a Mariano Puertas, Jefe de Programa de Planificación Urbana.
\item \textbf{Decisión Administrativa 366/22} (03/05/2022): designa al Dr. José Javier Haro, Director General de Tierras Fiscales y Bienes del Estado.
\item \textbf{Decisión Administrativa 19/23} (26/12/2023): con el nuevo mandato, vuelve a designar a Haro y a Puertas en sus cargos y nombra, entre otros, a Luis Eduardo Cornejo Revilla Córdoba como Director General de Inmuebles.
\item \textbf{Decreto 64/23} (10/12/2023; B.O.\ Nº 21.610 del 12/12/2023): designa a la \textbf{Dra. María Silvina Abilés}, Secretaria de Tierras y Bienes del Estado. Lo tramita el \textbf{Ministerio de Infraestructura} y lo refrendan su ministro y la Secretaria General de la Gobernación. Los Decretos 197/24 y 198/24 la designan interventora \textit{ad honorem} del Programa de Vivienda Popular Sociedad del Estado y la autorizan, junto con la Subsecretaria de Regularización Dominial Ing. Alicia Virginia Galli, a suscribir actos jurídicos en nombre de la Provincia.
\end{itemize}

El sitio oficial de la Municipalidad de La Caldera documenta la relación de trabajo entre el municipio y la Secretaría: consigna que el intendente Diego Sumbay\index{Sumbay, Diego Ramiro}, acompañado por el responsable de Medio Ambiente, Tierra y Hábitat municipal, Rodrigo Páez, se reunió con la secretaria de Tierras y Bienes del Estado; y que ambos organismos iniciaron la primera etapa del proceso de escrituración de terrenos municipales para una quincena de familias caldereñas.

La grafía oficial del apellido es \textbf{Abilés}, la misma que emplea el registro de la Cámara de Senadores; ``Avilés'', que aparece en parte de la cobertura periodística, es error de prensa y no una persona distinta.

Que la titular del organismo provincial de tierra fiscal trabaje con el municipio del departamento que representó ocho años en la Legislatura no es irregular ni sorprendente: es su función, y la regularización dominial de familias caldereñas es exactamente lo que el Plan ``Mi Lote'' se propone. Lo que el dato agrega es densidad, no imputación. La tercera posición --- la copropiedad del loteo Solanas --- sigue siendo, como estableció el capítulo \ref{cap:loteo}, una atribución periodística sin respaldo registral.

\section{Los tres procedimientos}

El Anexo I de la Resolución 9/16 ---publicado, como el de la 8/16, sólo como imagen escaneada vinculada al acto en el Boletín Oficial Nº 19.906, del 17 de noviembre de 2016--- establece que los procedimientos de asignación son tres, y los acota con cupos.

\begin{ficha}{Cupos y garantías --- Anexo I, Resolución 9/16}
\begin{itemize}[leftmargin=*,nosep]
\item \textbf{Casos generales} (art. 5º): sorteo público. Cupo \textbf{mínimo del 70\%} de los lotes de la urbanización.
\item \textbf{Casos especiales} (art. 6º): sorteo entre preseleccionados por circunstancias socioeconómicas comprobadas o antigüedad, con informe del Área Social y resolución fundada. Cupo \textbf{máximo del 28\%}.
\item \textbf{Asignación directa} (art. 7º): reservada a situaciones de vulnerabilidad que ``signifiquen riesgo a la integridad física, psíquica, salud y/o seguridad personal o de su grupo familiar'', acreditadas por informe socioambiental fundado. Cupo \textbf{máximo del 2\%}. Padrón exhibido cinco días para tachas. Se dispone por Resolución de la Subsecretaría con visto bueno del Secretario.
\item \textbf{Garantías del sorteo} (arts. 3º y 4º): escribano de Escribanía de Gobierno obligatorio, sin cuya presencia el acto no puede comenzar \textit{bajo pena de nulidad}; lugar de acceso público; padrones provisorios y definitivos exhibidos en sede y publicados en la web provincial y en el Boletín Oficial; domicilio constituido para observaciones e impugnaciones de cualquier interesado; once contenidos obligatorios del aviso.
\item \textbf{Obligaciones} (arts. 11 a 16): construir en un plazo no mayor a tres meses; prohibición de vender, ceder, alquilar o dar en comodato bajo pena de inhabilitación definitiva; visitas de constatación del Área Social; revocación por acto fundado con derecho de defensa.
\end{itemize}
\end{ficha}

El Anexo II es el modelo de Acta de Tenencia Precaria, y conviene mirarlo porque es el único instrumento del régimen que el beneficiario firma. Lo suscriben el Subsecretario de Tierra y Hábitat y el beneficiario; obliga a destinar el lote «única y exclusivamente al uso habitación» del grupo familiar declarado; deja al organismo las obras hechas si la tenencia se revoca, sin derecho a repetir gasto alguno; y da cinco días para restituir el lote libre de ocupantes desde el acto que ordene la revocación. Con un detalle que importa: \textbf{la cláusula primera, la que fija en cuántos meses hay que ocupar el lote, está en blanco en el modelo}, mientras el artículo 11º del Anexo I fija tres meses para empezar a construir. Cuál de los dos plazos rige lo decide, caso por caso, quien completa el acta.

Y el artículo 4º agrega dos plazos que el resto del régimen no tiene: la Subsecretaría debe pedir el escribano a Escribanía de Gobierno con \textbf{siete días} de anticipación, y el administrador del sorteo debe remitir los resultados con los padrones de beneficiarios dentro de los \textbf{dos días} posteriores. Son los únicos plazos en días que el procedimiento fija por sí mismo, y los dos rigen hacia adentro de la administración; el que corre para el vecino ---cuántos días se exhibe el padrón--- sigue remitido al Anexo de la Resolución 82/16.

Conviene distinguir dos institutos cuyos nombres se parecen y que este relevamiento confundió antes de leer el anexo. La \textbf{asignación directa} del artículo 7º --- 2\% de cupo, riesgo vital, tenencia precaria sobre lote de urbanización nueva --- no es lo mismo que la \textbf{adjudicación en venta por adjudicación directa} del artículo 1º inciso a de la Ley 2616, que es regularización dominial de ocupantes preexistentes y termina en título de propiedad.

\section{El baremo}

El Anexo I de la Resolución 8/16 fija las tipologías de puntaje. Situación habitacional: vivienda precaria 30 puntos, vivienda con baño o cocina compartidos 20, vivienda deficitaria --- baño fuera, sin arrastre de agua, agua fuera, sin luz permanente --- 30; tenencia en situación de calle 50, inquilino 20, ocupante en relación de dependencia 10; habitabilidad crítica --- tres o más personas en el mismo cuarto --- 50. Grupo familiar: 15 por cada menor de 18 a cargo y 15 por cada adulto mayor de 70. Enfermedad grave o discapacidad: 15 por cada integrante en cada supuesto. Antigüedad de inscripción: 0,5 por cada mes cumplido.

El baremo tiene un hueco que conecta los dos extremos del libro.

La tipología de tenencia puntúa tres situaciones y \textbf{no puntúa la ocupación de hecho}. El formulario de inscripción pregunta expresamente si el solicitante es ocupante gratuito con permiso u ocupante de hecho sin permiso; ninguna de las dos recibe puntaje alguno.

Y el ocupante de hecho es, precisamente, el primer orden de preferencia del artículo 2º de la Ley 2616 y el fundamento expreso con el que el Decreto 63/25 adjudicó el barrio Juan Manuel de Rosas. La provincia tiene dos vías que se ignoran mutuamente: por la del sorteo con puntaje, ocupar un terreno fiscal no suma nada; por la de la ley de 1951, ocupar es lo que encabeza la preferencia. Quien ocupa no compite en la lista: espera a que llegue un decreto.

Tres observaciones menores. La antigüedad no tiene tope --- treinta años acumulan 180 puntos, cifra comparable al máximo por necesidad --- y la Resolución 7/16 obliga a actualizar o ratificar el legajo todos los años y dispone la baja automática, con pérdida de toda la antigüedad, si pasan tres años sin hacerlo. La tabla de características de la vivienda enumera los ítems A, C, D y E: no hay ítem B. Y el formulario pregunta si la vivienda está a menos de tres cuadras de un basural, si estuvo \textbf{en zona inundable en los últimos doce meses} y cuánto gana el solicitante; el baremo no puntúa ninguna de las tres.

En un departamento cuyo conflicto central es la inundación de barrios enteros, y donde el Estado fue condenado por el manejo de una microcuenca, el sistema de asignación de tierra fiscal registra la exposición al riesgo hídrico y le asigna cero.

\section{Las incompatibilidades y el requisito de residencia}

El artículo 2º del Anexo I de la Resolución 7/16, inciso d, establece que no podrán ser adjudicatarios de lotes de interés social, mientras ejerzan sus cargos, el Gobernador, el Vicegobernador, los Ministros, Secretarios y Subsecretarios de Estado, Directores Generales, Directores, Jefes de Programa y toda persona con designación vigente en carácter de Autoridad Superior del Poder Ejecutivo provincial, municipal o nacional, los presidentes y directores de empresas y sociedades del Estado y de entidades autárquicas, \textbf{sus cónyuges o convivientes y su grupo familiar declarado}. El mismo artículo excluye, en su inciso c, a quienes ---titular o grupo familiar--- estén incursos en contravenciones o infracciones vinculadas con la \textbf{ocupación ilegal} de inmuebles públicos o privados, y admite dos excepciones a la regla de no tener inmuebles: quien transfirió gratuitamente sus derechos tras un divorcio, y quien es titular de hasta un tercio de un inmueble que razonablemente no puede usarse como vivienda.

Los artículos 4º inciso e y 6º exigen Certificado original de Residencia y Convivencia expedido por la Seccional Policial, emitido con un máximo de treinta días de anticipación, y el formulario registra la residencia en la localidad donde se solicita el lote y desde qué fecha.

Y el artículo 8º agrega tres reglas que importan para este libro. El trámite es personal y gratuito, y se hace \textbf{en la Subsecretaría de Tierra y Hábitat o en el municipio donde se lleve adelante la urbanización, «en caso de Convenio vigente»} (inc.\ a): el municipio puede ser la ventanilla del circuito, de modo que los convenios que el apéndice \ref{cap:pedidos} pide a la Secretaría dirían si La Caldera o Vaqueros lo fueron alguna vez. El acceso es \textbf{únicamente por sorteo público, salvo los casos de asignación directa} (inc.\ j). Y el artículo 9º remite las impugnaciones a la Ley 5.348.

\begin{cautela}
El inciso d prohíbe que una autoridad superior \textit{reciba} un lote de interés social; no dice nada sobre que una autoridad sea propietaria de suelo privado. La situación descripta en el capítulo \ref{cap:politica} no está alcanzada por este artículo. Lo que el artículo sí establece, y vale como criterio, es que el propio organismo consideró necesario escribir que quien ejerce autoridad sobre la distribución de tierra no puede ser destinatario de ella.

Sobre la residencia corresponde una corrección en favor de los funcionarios cuestionados en el apartado siguiente. La Resolución 7/16 \textbf{no contiene} una prohibición de adjudicar a habitantes de otros municipios, como se afirmó públicamente. Contiene una exigencia de acreditación policial de residencia y un registro obligatorio de su antigüedad. No hay, al menos en esta norma, un impedimento legal que alguien haya violado. \textit{(El apartado «El anexo que faltaba», más abajo, matiza esto: el reglamento marco de la Resolución 82/16 sí exige cuatro años de residencia ininterrumpida en la localidad donde se solicita el lote, aunque sólo para los sorteos. La conclusión no cambia, y la advertencia corresponde hacerla también acá.)} Y sus anexos, como los de las Resoluciones 8/16 y 9/16, están en el repositorio del Boletín sólo como imagen escaneada: el formulario de inscripción y la lista de requisitos existen publicados, pero ningún buscador los encuentra.
\end{cautela}

\section{Juan Manuel de Rosas}

En abril de 2025, días antes de la elección provincial, se realizó un acto de firma de decretos de adjudicación de terrenos en el barrio Juan Manuel de Rosas, en la zona norte de la ciudad de Salta, contiguo al río Vaqueros. La firma estuvo a cargo de la secretaria de Tierras y Bienes del Estado; el discurso de entrega, a cargo del entonces senador provincial por La Caldera y candidato a renovar la banca, Miguel Calabró, quien destacó que entre los beneficiarios había vecinos vaquereños y caldereños residentes en ese barrio de la capital.

La presidenta del centro vecinal del barrio, María Fernanda Mamaní, cuestionó públicamente la atribución del logro, sostuvo que lleva más de siete años gestionando la regularización, que en el acto se firmaron decretos para treinta vecinos y que \textbf{veintiséis de esos treinta son vecinos de La Caldera que reciben un terreno en Salta capital}, y afirmó que existe una norma que prohíbe adjudicar a habitantes de otros municipios y que está prohibido realizar actos de este tipo dentro de los treinta y cinco días previos a una elección.

\begin{ficha}{Decreto Nº 63/25}
\begin{itemize}[leftmargin=*,nosep]
\item \textbf{Fecha:} 6 de febrero de 2025. Boletín Oficial Nº 21.894 del 13/02/2025. Expediente Nº 1100388-113156/2024-0 y agregados. El anexo se publica como archivo vinculado al decreto en la edición web: cuatro folios de una copia certificada escaneada ---con el sello «Es copia» de la Dirección General de Leyes y Decretos---, que a diferencia de los anexos de las Resoluciones 7, 8, 9 y 82 de 2016 trae una capa de texto de reconocimiento óptico.
\item \textbf{Objeto:} adjudica en venta inmuebles de propiedad de la Provincia en el Barrio Juan Manuel de Rosas, departamento Capital.
\item \textbf{Marco:} el art.\ 37 de la Constitución provincial y el Plan «Mi Lote» del Decreto 399/20, que tiene entre sus objetivos la regularización dominial de quienes ocupan sin título.
\item \textbf{Fundamento de la vía:} la Dirección General Unidad de Regularización Dominial informó que los solicitantes reúnen los requisitos \textbf{en los términos del art. 2º inc. 1 de la Ley 2616, por revestir el carácter de actuales ocupantes}. Según informes de la Dirección General de Inmuebles y del Instituto Provincial de Vivienda, \textbf{no son titulares ni beneficiarios de inmuebles u otras soluciones habitacionales}.
\item \textbf{Condiciones:} hasta cien cuotas mensuales; obligación de habitar con el grupo familiar en forma continua y permanente como vivienda única; prohibición de enajenar, locar o dar en comodato sin autorización hasta la cancelación de la hipoteca; deudas e impuestos a cargo de los adjudicatarios desde la notificación; escrituras por Escribanía de Gobierno, con hipoteca exenta de impuestos y las obligaciones transcriptas; caducidad automática ante cualquier incumplimiento, con reversión del inmueble \textbf{y de todo lo construido o plantado}, sin derecho a indemnización.
\item \textbf{Firmantes:} Sáenz --- Camacho, ministro de Infraestructura --- López Morillo, secretaria general de la Gobernación.
\item \textbf{Modificado por:} Resolución Delegada Nº 567D/26 (B.O.\ 22.224, 13/07/2026), que subsana errores materiales en los anexos de dieciocho instrumentos, entre ellos éste. Su contenido se analiza más abajo.
\end{itemize}
\end{ficha}

\begin{cautela}
\textbf{Una discordancia que conviene poner primero.} El acto que describe la prensa es de abril de 2025 y habla de decretos firmados ese día para treinta vecinos. El Decreto 63/25 es del 6 de febrero y comprende ciento tres parcelas. No son, en sentido estricto, el mismo acto: el de abril pudo ser la entrega de decretos dictados antes ---éste entre ellos--- o la firma de otros que este libro no encontró. La regularización del barrio, por lo demás, avanza por tandas: la Secretaría entregó en 2023 un plano con más de ochocientos catastros individuales, y en julio y noviembre de 2025 hubo nuevas entregas y firmas de escrituras, de sesenta y de setenta y dos familias. El Decreto 63/25 es, entonces, uno de varios instrumentos del mismo proceso, y la primera explicación es la más probable. Lo que sigue usa el Decreto 63/25 porque es el único instrumento de adjudicación del barrio que se obtuvo; verificar la cifra de veintiséis sobre treinta requiere identificar los decretos del acto de abril.
\end{cautela}

El anexo nominal comprende \textbf{ciento tres parcelas} en las manzanas 430 a 475, con matrículas entre la 188.191 y la 188.612. Sus columnas son manzana, parcela, matrícula, expediente, titular, DNI, condición porcentual, co-titular, DNI y valor.

\begin{cautela}
\textbf{El anexo no consigna domicilio.} No hay columna de dirección, localidad ni municipio. El documento que este libro fue a buscar para verificar la afirmación de los veintiséis de treinta no permite verificarla ni desmentirla.

Y una decisión editorial que se explicita: el anexo identifica con nombre completo y documento a un centenar de personas de bajos ingresos. Algunos apellidos coinciden con los que la nota periodística mencionó. Este libro no publica esa lista, no la coteja por apellido y no nombra a ningún adjudicatario. Sin el dato de domicilio, un cotejo onomástico produce coincidencia y no identificación; y producir una lista de personas señaladas como beneficiarias irregulares, en la posición social más vulnerable del expediente y sin capacidad de defensa, sería atribuir a individuos concretos una imputación que ni siquiera está probada respecto de los funcionarios.

Lo que efectivamente demuestra el anexo, leído en agregado, es otra cosa. Los expedientes tienen dos series. Noventa y una parcelas tramitan sólo por expedientes de la serie actual, iniciados entre 2021 y 2024, y casi la mitad del total es de 2022. Doce parcelas arrastran además trece expedientes de la serie anterior, con fechas de 1995, 2000, 2001, 2004, 2005, 2006 (dos), 2007 (dos), 2014, 2015 (dos) y 2018; una de ellas acumula tres expedientes, de 1995, 2014 y 2024. Un expediente iniciado en 1995 y resuelto en 2025 significa treinta años para regularizar a una familia que ocupaba el terreno. \textit{(El escaneo del anexo tiene números de expediente parcialmente ilegibles: los recuentos por año son aproximados; los de la serie anterior, no.)}

Eso obliga a leer el fundamento con cuidado. Que la mayoría de los expedientes sea de 2022 en adelante no indica que la ocupación sea reciente: coincide con el censo que la Secretaría inició ese año para regularizar un barrio que ya tenía treinta años. La demora larga está documentada en doce parcelas, y en una de ellas el primer expediente es de 1995. Alcanza para desactivar la lectura de que el acto habría sido una operación electoral improvisada; no alcanza para afirmar que todos los casos arrastran tres décadas de trámite.

Lo que la demora sí documenta, y es lo que el libro retiene, es el tiempo de respuesta del circuito: se mide en décadas. Ése es el contraste con el otro, donde una parcela de El Durazno pasa de monte a escritura en el lapso de un ciclo comercial.
\end{cautela}

Los valores oscilan aproximadamente entre 2,8 y 7,9 millones de pesos por parcela, con hipoteca en garantía del saldo, determinados con el coeficiente de disminución del treinta por ciento de la Ley 2616. Las condiciones de titularidad registradas son del 100 por ciento, del 50 y 50, del 33,33 entre tres y del 25 entre cuatro.

Queda una cuestión normativa, que el anexo no permite contrastar y que este libro no intenta contrastar con los datos de los adjudicatarios. El artículo 2º de la Ley 2616, de 1951 y vigente según el Anexo I del Digesto Jurídico, coloca en segundo lugar de preferencia a ``los argentinos nativos'': un criterio que hoy sería difícil de sostener frente al artículo 20 de la Constitución Nacional.

\section{Lo que la Resolución 567D/26 corrige}

Se obtuvo el texto de la resolución que subsana los errores materiales del anexo del Decreto 63/25 y de otros diecisiete instrumentos.

\begin{ficha}{Resolución Delegada Nº 567 D --- B.O.\ Nº 22.224, 13 de julio de 2026}
Dictada el 6 de julio de 2026 por el Jefe de Gabinete de Ministros, expediente 1100388-214146/2024-0, con firma de \textbf{Camacho}.

\textbf{Origen:} la \textbf{Secretaría de Tierra y Bienes del Estado} analizó los decretos de adjudicación en venta de lotes fiscales de interés social y de caducidad de adjudicaciones, y \textbf{detectó errores materiales en la confección de sus Anexos}.

\textbf{Fundamento:} artículo 76 de la Ley 5.348/78 de Procedimientos Administrativos, que permite rectificar de oficio, en cualquier momento, los errores materiales y aritméticos \textbf{siempre que la enmienda no altere lo sustancial del acto}. La resolución los califica de ``errores materiales no esenciales''. Facultades delegadas por Decreto 41/96, modificado por sus similares 1.105/02 y 1.575/08.

\textbf{Alcance:} dieciocho instrumentos ---con veinte rectificaciones, porque dos de ellos llevan más de una---, con fechas que van de \textbf{1991 a 2025}: Decretos 871/91, 521/99, 2067/05, 2068/05, 2683/10, 1976/11, 3354/12, 1608/14, 3824/14, 1359/18, 992/19, 1699/19, 373/21, 218/23, 233/23, 428/23 y 63/25, más la Resolución Delegada 156D/23. Dos anexos: Capital e Interior, \textbf{publicados como imagen escaneada vinculada al acto} y certificados como copia por el director de Leyes y Decretos.
\end{ficha}

\begin{cautela}
Nota de método, que corresponde hacer porque este libro se equivocó y conviene decir cómo.

El detalle de las rectificaciones llegó primero a este relevamiento a través de un \textbf{reporte de análisis} elaborado sobre la resolución, no del anexo original. Sobre esa base se escribió un primer borrador de esta sección. Obtenido después el \textbf{Anexo publicado} --- documento SA100053553 del Boletín Oficial ---, el cotejo mostró que el reporte \textbf{atribuía cambios de parcela y de matrícula que el anexo no contiene}.

Concretamente: el reporte afirmaba que en Juan Manuel de Rosas la parcela pasaba de 8 a 18, que en Congreso Nacional la matrícula pasaba de 18.827 a 8.827 y que en Coronel Juan Solá la parcela pasaba de 5 a 15. En el anexo, las tres columnas figuran \textbf{idénticas en las líneas ``dice'' y ``debe decir''}: lo único que cambia en esos tres casos es la grafía de un nombre.

Se deja constancia porque es exactamente el riesgo que la advertencia de procedencia anticipaba, y porque un documento intermedio que resume no es una fuente. Lo que sigue está tomado del anexo.
\end{cautela}

\subsection*{Qué se corrigió en Juan Manuel de Rosas}

De las ciento tres parcelas del Decreto 63/25, el Anexo I registra \textbf{dos} rectificaciones. Las dos corrigen la grafía del nombre de pila del titular: en una faltaba una letra; en la otra, el nombre estaba escrito con otra ortografía. Este libro no reproduce los nombres ni los datos de las parcelas, por la razón que se explicó más arriba. \textbf{En los dos casos, manzana, parcela, matrícula, expediente, documento, condición y valor permanecen sin cambio.}

Son correcciones ortográficas y nada más. El decreto que este capítulo analiza no fue alterado en ninguno de sus datos sustanciales.

\subsection*{Qué sí cambió de inmueble}

Dos casos del anexo modifican la identificación del bien, y ninguno es de La Caldera ni de Juan Manuel de Rosas.

\begin{ficha}{Barrio Leopoldo Lugones --- Decreto Nº 218/23}
\textbf{Dice:} manzana 364 C, \textbf{parcela 4}, \textbf{matrícula 90.279}, expediente 131-39395/04.

\textbf{Debe decir:} manzana 364 C, \textbf{parcela 3}, \textbf{matrícula 187.052}, expedientes 131-39395/04 y 388-47308/23.

Titular y co-titular, documentos, condición porcentual y valor (\$1.464.050) permanecen idénticos.
\end{ficha}

\begin{ficha}{Barrio Los Profesionales --- Decreto Nº 521/99}
\textbf{Dice:} manzana 90A, \textbf{parcela 1B}, \textbf{matrícula 123.194}, expediente 131-11542/97.

\textbf{Debe decir:} manzana 90A, \textbf{parcela 15}, \textbf{matrícula 186.509}, mismo expediente.

Titulares, documentos, condición y valor (\$2.829) permanecen idénticos.
\end{ficha}

\begin{cautela}
Aquí sí el objeto cambia. Entre la sanción del decreto en 2023 y la rectificación en julio de 2026, el instrumento consignó una parcela y una matrícula distintas de las que corresponden. La enmienda se practicó por la vía del artículo 76 de la Ley 5.348/78, que la admite \textbf{siempre que no altere lo sustancial del acto}.

Este libro no discute la validez jurídica de esa calificación, que no le compete, y admite que el sentido más plausible es benigno: que el expediente siempre acreditó la parcela 3 y que el anexo del decreto la transcribió mal. La aparición de un segundo expediente acumulado en la línea corregida apunta en esa dirección.

Registra, en cambio, un rasgo del sistema que este capítulo viene describiendo: \textbf{el número de parcela y la matrícula pueden corregirse sin rehacer el decreto de adjudicación}. Y registra que la detección tardó tres años en un caso, veintisiete en el otro, y hasta treinta y cinco en los más viejos del acto.
\end{cautela}

\subsection*{Y qué revela el conjunto}

Con el anexo a la vista, la mayoría de las veinte rectificaciones son de grafía de nombres y apellidos: \textit{Telles} por \textit{Tellez}, \textit{Reynoso} por \textit{Reinoso}, \textit{Nacy} por \textit{Nancy}, \textit{Cruz} por \textit{Crus}, \textit{Mirian} por \textit{Miriam}, \textit{Alerto} por \textit{Alberto}, \textit{Exequiel} por \textit{Ezequiel}, \textit{Irma} por \textit{Yrma}. Junto a ellas, cinco tipos que sí tienen peso:

\begin{itemize}[leftmargin=*]
\item \textbf{Dos cambios de parcela y matrícula:} Leopoldo Lugones y Los Profesionales, ya vistos.
\item \textbf{Una conversión del precio en «sin cargo»:} en el Barrio Martel de Tartagal (Decretos 2067/05 y 3354/12), el valor de \$8.617,27 pasa a «sin cargo», junto con la corrección del apellido y la incorporación de un expediente de 2024.
\item \textbf{Dos correcciones de documento de identidad:} en Finca La Floresta cambia el número de documento de la titular; en Santa Rita Sur, el de la co-titular.
\item \textbf{Una supresión de co-titular:} en Barrio Convivencia se elimina a un co-titular que figuraba junto a una titular que ya tenía el cien por ciento de la condición.
\item \textbf{Incorporación de expedientes acumulados:} en La Loma se agrega un expediente de 2024 donde el decreto de 1991 no consignaba ninguno ---su valor sigue expresado en australes---; en Siglo XXI y en Martel se agregan expedientes de 2019 y 2024.
\end{itemize}

Las dos correcciones de documento cambian el dato con que el Estado identifica a quien adjudica. Pueden deberse a un documento renovado o a un error de carga, y el anexo no dice cuál. Lo que el dato muestra es que \textbf{corregir ese registro llevó, en un caso, desde 2011 hasta 2026}.

Y una supresión de co-titular tampoco es cosa menor. En Barrio Convivencia figuraba una segunda persona junto a la titular durante doce años, en un decreto de 2014, y en 2026 se la eliminó. Quién era y por qué figuraba está en el expediente 131-44942/06.

\begin{cautela}
Dos advertencias sobre el alcance de este cotejo.

\textbf{La primera:} el anexo es una imagen escaneada, y en la versión que este libro usó primero varias tablas ---17 de Octubre, Los Profesionales, Siglo XXI, Virgen del Rosario y el segundo caso del Barrio Martel de Tartagal--- no resultaron legibles. Una lectura posterior del mismo documento las recuperó: tres son de grafía u orden de nombres, Los Profesionales cambia parcela y matrícula, y en el segundo caso de Martel \textbf{el valor pasa de \$8.617,27 a «sin cargo»}. El dato que el reporte intermedio anticipaba, y que este libro se había negado a afirmar sin el documento, resultó cierto. \textbf{Es la única rectificación del acto que mueve dinero}, y se practicó por la vía de un artículo que sólo admite enmiendas que no alteren lo sustancial. Una discordancia menor: el anexo titula «Res.\ 159D/23» el caso de Finca La Floresta, y el texto de la resolución enumera la «Resolución Delegada 156D/23».

\textbf{La segunda:} ninguna de las rectificaciones legibles corresponde al departamento La Caldera. Los anexos cubren Capital, Cerrillos, General San Martín y Rivadavia. El departamento aparece sólo de manera indirecta, a través del Barrio Juan Manuel de Rosas --- que está en la Capital y donde, según el centro vecinal, veintiséis de treinta adjudicatarios serían vecinos caldereños. Esa afirmación sigue sin poder verificarse: el anexo rectificado tampoco consigna domicilio.
\end{cautela}

Queda, sí, una continuidad de firma. \textbf{Camacho} firma esta resolución como Jefe de Gabinete en julio de 2026. Es quien contrafirmó el Decreto 63/25 en febrero de 2025, cuyo anexo aquí se corrige; quien firmó, como Ministro de Infraestructura, la Resolución 646D/21 que aprobó el plano del loteo El Durazno (capítulo \ref{cap:loteo}); y quien convocó, como Jefe de Gabinete, la reunión de enero de 2026 en la que se aceptó el plan de represas de retardo propuesto por el intendente de La Caldera.

Tres actos y una reunión, dos cargos, un mismo funcionario. Como en el caso de la funcionaria del área hídrica, se consigna sin connotación: describe dónde está concentrada la decisión sobre suelo en la provincia, no la conducta de quien decide.


\section{Y en 1943 fue a juicio para prescribir}

El capítulo documenta que en 2015 la Provincia declaró operada a su favor la prescripción
adquisitiva del terreno de una escuela emplazada desde 1968 sobre suelo ajeno. No fue la
primera vez.

\begin{ficha}{Edicto de posesión treintañal --- B.O.\ Nº 1.994, 19 de marzo de 1943}
El \textbf{Gobierno de la Provincia de Salta}, representado por el doctor Raúl Fiore Moules,
deduce juicio de posesión treintañal ante el Juzgado de 2ª Nominación en lo Civil, a cargo del
doctor Ricardo Reimundín, sobre el inmueble del pueblo de La Caldera «donde se encuentra
emplazado el edificio ocupado por la \textbf{Comisaría de Policía}».

\textbf{Dimensiones:} cuarenta metros de frente por doscientos de fondo --- \textbf{ocho mil
metros cuadrados}.

\textbf{Límites:} «Norte, propiedad de David Villa Gómez; Sud y Este, propiedad de
\textbf{Daniel Linares\index{Linares, Daniel}}; Oeste, calle principal del pueblo de La Caldera.»

Publicación: treinta días en dos diarios y \textbf{una sola vez} en el Boletín Oficial.
\end{ficha}

\textbf{La primera: el Estado ocupaba sin título su propia comisaría.} Ocho mil metros
cuadrados sobre la calle principal del pueblo, ocupados durante al menos treinta años ---
es decir, desde 1913 o antes --- sin instrumento que lo acreditara. Y probablemente desde
antes: en agosto de 1910 la Provincia aprobó el contrato con Moisés Vera para construir
«el edificio para la comisaría de La Caldera» (B.O.\ Nº 178), sin decir sobre qué terreno. La misma situación que
setenta y dos años después se resolverá para una escuela.

\textbf{La segunda: en 1943 tuvo que ir a juicio.} A diferencia de 2015, cuando un decreto
del Poder Ejecutivo bastó como título por aplicación de la Ley nacional 21.477, aquí la
Provincia \textbf{demanda ante un juez civil}, con abogado apoderado, edictos por treinta
días y traslado a quien pudiera oponerse.

El régimen cambió en 1977. Pero conviene ver qué significa: \textbf{en 1943 el Estado se
sometía al mismo procedimiento que un particular para adquirir por prescripción, y desde
1977 no}. La asimetría que este capítulo describe entre el ocupante particular --- que debe
litigar --- y el Estado --- que dicta un decreto --- tiene fecha de nacimiento.

\textbf{Y la tercera: el lindero, por dos costados, es Daniel Linares\index{Linares, Daniel}.} El presidente de la
comisión de defensas del río de 1910, miembro del cuerpo municipal en 1922, propietario de
Los Porongos en 1932 y de El Angosto en 1938, es en 1943 dueño de los terrenos que rodean a
la comisaría del pueblo. Su trayectoria, en el apéndice \ref{cap:personas}, cubre treinta y tres años.

\subsection*{Y lo que costó}

El archivo permite algo que casi nunca permite: saber cuánto salió un trámite. Dos decretos
de 1943 liquidan los gastos de ese juicio de posesión treintañal.

El \textbf{Decreto 7223-H} del 15 de abril paga \textbf{setenta pesos} a la administración del
diario \textit{La Provincia} por la publicación del edicto, imputados a la cuenta «deudores
juicios varios». El \textbf{Decreto 7310-H} del 14 de mayo paga \textbf{doce pesos} a
\textbf{Teófilo Reyes\index{Reyes, Teófilo}}, Juez de Paz Propietario de La Caldera, por diligenciar un oficio
librado en esa causa.

Ochenta y dos pesos, más el trabajo del apoderado. Póngase esa cifra al lado de las otras que
este libro tiene para el período: la renta anual del municipio rondaba los dos mil
ochocientos treinta y cinco pesos en 1937, y las obras de defensa del pueblo contra el río
costaron setecientos en 1935; las de La Calderilla de 1938, \textbf{doscientos cincuenta pesos por cincuenta metros}: cinco pesos el metro.

\textbf{Regularizar el título de ocho mil metros cuadrados en la calle principal del pueblo le costó al Estado provincial lo que unos dieciséis metros de defensa sobre el río}, al precio de la de 1938. Y aun así, treinta
y cuatro años después una ley nacional lo relevó de tener que hacerlo por vía judicial.

Conviene retener además el nombre del juez de paz que diligenció el oficio, porque el mismo
archivo lo registra en otro papel: \textbf{Teófilo Reyes\index{Reyes, Teófilo}} comparte apellido con el Teodoro
Reyes que en 1909 figuraba como comisario auxiliar del partido de Los Porongos. Coincidencia
de apellido, como todas las del apéndice \ref{cap:personas}.

\section{Cómo compraba tierra el Estado en 1935}

Este capítulo describe cómo circula la tierra fiscal ---el sorteo y la regularización de ocupantes--- y cómo el Estado adquiere la que ocupa ---la prescripción administrativa a su favor---. Hubo otra vía de adquisición, y en
1935 se la usaba de manera rutinaria.

\begin{ficha}{Licitación de la Dirección General de Obras Públicas --- B.O.\ Nº 1.590, 28 de junio de 1935}
«Se avisa a los interesados en presentar propuestas que está abierta la licitación para la
\textbf{adquisición de los terrenos} destinados a \textbf{edificación escolar} en las
siguientes localidades»: siguen unas treinta, entre ellas \textbf{Vaqueros}.

Y una segunda lista, para terrenos destinados a \textbf{asistencia social}, en la que figura
\textbf{La Caldera}.

Apertura de propuestas: lunes 15 de julio de 1935, en la Dirección General de Obras Públicas.
\end{ficha}

\textbf{El Estado compraba el suelo que necesitaba, y lo compraba por licitación pública.} No
expropiaba ni prescribía: publicaba un aviso, abría propuestas y adquiría. Y lo hacía en
escala provincial, con una sola licitación para treinta localidades.

Ese procedimiento conviene tenerlo presente al leer el resto de este capítulo, porque hace de
contraste con los tres circuitos contemporáneos.

En 1935 el Estado \textbf{llama a licitación para comprar} terrenos escolares. En 2015, para regularizar el de
otra escuela emplazada desde 1968 sobre suelo ajeno, \textbf{declara operada a su favor la
prescripción adquisitiva} por decreto que vale como título. Y en el medio, la Ley 4486 de
1972 \textbf{expropia} seiscientas ochenta y seis hectáreas con toma de posesión previa a la
discusión del precio.

Comprar, expropiar, prescribir: \textbf{tres formas de que el Estado obtenga suelo, en orden
decreciente de lo que le cuesta y de lo que el propietario puede discutir.} Este libro no
sostiene que haya habido un desplazamiento deliberado de una hacia otra --- son institutos
distintos para situaciones distintas, y los tres son legales --- pero registra que la primera
es la que el archivo muestra en los años treinta y la última la que muestra en el presente.

\section{El perilago ya se cedía en 1980}

Este capítulo documenta que en 2015 la Provincia cedió en comodato gratuito cuarenta y una
hectáreas fiscales por veinte años para una cancha de golf. La figura --- ceder tierra fiscal
del perilago a una entidad privada --- es treinta y cinco años más vieja.

\begin{ficha}{Decreto Nº 1591 del 10 de noviembre de 1980}
Expediente 11-16514/80. Se otorga \textbf{permiso precario de uso} al «Club de Regatas Campo
Alegre», con personería jurídica reconocida el año anterior, sobre una fracción de
\textbf{una hectárea y media} del catastro 1575 del departamento La Caldera, \textbf{sobre el
perilago del Dique Ing.\ Alfonso Peralta}, según plano 148 inserto a fojas 5.

Carácter: \textbf{revocable}, con plazo máximo de \textbf{cinco años} desde la aceptación
expresa de la beneficiaria.
\end{ficha}

Conviene poner los dos actos uno al lado del otro, porque el instrumento es el mismo y las
condiciones no.

\needspace{6\baselineskip}
{\small
\begin{longtable}{@{}p{3.0cm}p{4.6cm}p{4.6cm}@{}}
\toprule
 & \textbf{1980} & \textbf{2015} \\
\midrule
\endhead
Beneficiario & Club de Regatas Campo Alegre & Proyecto 1 S.A. \\
Superficie & 1,5 ha & 41 ha \\
Plazo & \textbf{5 años} & \textbf{20 años} \\
Carácter & \textbf{precario y revocable} & comodato gratuito \\
Destino & no consta en el acto & cancha de golf \\
\bottomrule
\end{longtable}
}

\textbf{Casi veintiocho veces más superficie y cuatro veces más plazo}, y la figura pasa de
permiso precario revocable a comodato.

Este libro no sostiene que una cosa lleve a la otra ni que haya continuidad deliberada entre
ambos actos: median treinta y cinco años, dos regímenes constitucionales y beneficiarios de
naturaleza distinta --- una asociación civil con personería y una sociedad anónima. Lo que
registra es que \textbf{la tierra del perilago se viene cediendo a privados desde que el
embalse existe}, y que el instrumento se volvió más generoso: más años, más hectáreas, menos
revocable.

\section{Cuarenta y una hectáreas fiscales para una cancha de golf}

Este capítulo describe dos circuitos de tierra fiscal: el sorteo y la regularización de ocupantes. Hay otra vía de salida de la tierra fiscal, y no está en el reglamento.

\begin{ficha}{Decreto Nº 2949 del 24 de agosto de 2015 --- B.O.\ Nº 19.614}
Otorga en \textbf{comodato} a la firma \textbf{Proyecto 1 S.A.}, CUIT 30-70916368-6, ejecutora del proyecto turístico ``Hotel La Caldera'', el inmueble de \textbf{matrícula 4.425 del departamento La Caldera}, \textbf{propiedad de la Provincia de Salta}.

\textbf{Superficie:} fracción A de 7 ha 1.101 m² y fracción B de 34 ha 7.157 m². \textbf{Total: 41 hectáreas 8.258 m²}, según plano de mensura y desmembramiento Nº 000943.

\textbf{Destino exclusivo:} construcción de una \textbf{cancha de golf de 18 hoyos} e instalaciones complementarias. \textbf{Plazo del comodato: veinte años.} Plazo para construir: cuatro años, bajo pena de restitución inmediata sin derecho a reclamo ni indemnización.

Fundamento: la Ley 6.064 ``prevé medidas promocionales para las empresas beneficiarias, como la locación a precio de fomento o \textbf{cesión en comodato de los bienes del dominio privado del Estado}''; y el artículo 22 del Régimen de Contabilidad faculta al Poder Ejecutivo a celebrar estos convenios cuando los bienes ``no tengan destino específico y la cesión redunde en beneficios para la comunidad''.

Antecedente: Resolución 62/01, que declaró de interés turístico el proyecto. Expedientes iniciados en 2010. Firman Urtubey, Parodi, Saravia, Ovejero y Simón Padrós.
\end{ficha}

\begin{cautela}
El comodato es, por definición del propio decreto, ``un préstamo de uso \textbf{gratuito}''.

De modo que en agosto de 2015 la Provincia entregó cuarenta y una hectáreas de tierra fiscal, sin cargo y por veinte años, a una sociedad anónima, para que construyera una cancha de golf. Seis meses antes, por el Decreto 463/15, esa misma sociedad había obtenido exención de tributos provinciales y créditos fiscales para ampliar su hotel.

\textbf{No se afirma aquí que ninguno de los dos actos sea irregular.} La Ley 6.064 prevé expresamente la cesión en comodato como medida promocional, el artículo 22 del régimen de contabilidad la autoriza, intervinieron el Fondo de Administración de Bienes, el Coordinador General de Tierra y Bienes y dos ministerios, y el decreto exige construir en cuatro años bajo pena de restitución. Es una política turística con marco legal, procedimiento y control de cumplimiento: el reglamento de la Ley 6.064 (Decreto 2448/85) prohíbe cambiar el destino del beneficio (art.\ 29), hace perderlo si los bienes se transfieren o arriendan (art.\ 18) y encarga a la Dirección Provincial de Turismo informes semestrales. Si la cancha se construyó en los cuatro años del decreto ---es decir, antes de agosto de 2019---, esos informes lo dicen.

Lo que corresponde es ponerlo al lado de lo que este mismo capítulo documenta sobre el otro circuito.

\begin{center}\small
\begin{tabular}{@{}p{4.0cm}p{4.0cm}p{4.0cm}@{}}
\toprule
 & \textbf{Lote de interés social} & \textbf{Comodato turístico} \\
\midrule
Superficie & Un lote urbano & \textbf{41,8 hectáreas} \\
Requisitos & Argentino nativo, cuatro años de residencia, sin inmuebles en el país ni en el exterior, sin haber recibido ayuda habitacional (Anexo de la Res.\ 82/16, analizado más abajo) & Ser la firma ejecutora de un proyecto declarado de interés turístico \\
Procedimiento & Sorteo con puntaje de vulnerabilidad, o regularización de ocupación & Solicitud y decreto \\
Tiempo & Décadas, según los casos documentados & Cinco años desde los expedientes de 2010; catorce desde la declaración de interés turístico \\
Título & Tenencia precaria, y adjudicación en venta posterior & Uso gratuito por veinte años \\
\bottomrule
\end{tabular}
\end{center}

No son institutos comparables en su naturaleza jurídica --- uno transfiere dominio, el otro cede uso --- y este libro no los equipara. Los pone en la misma página porque \textbf{provienen del mismo patrimonio}: la tierra del Estado provincial en el departamento La Caldera. Y porque la tesis de este trabajo es exactamente ésa: que la intensidad regulatoria de cada circuito de suelo es inversamente proporcional a la capacidad de sus beneficiarios.
\end{cautela}

\begin{cautela}
Dos precisiones que el decreto permite y conviene no pasar por alto.

\textbf{La primera:} el considerando dice que el Poder Ejecutivo puede ceder bienes ``\textbf{cuando no tengan destino específico} y la cesión redunde en beneficios para la comunidad''. Que esas cuarenta y una hectáreas fiscales no tuvieran destino específico en 2015 es una afirmación de la administración, no un hecho verificado por este libro. Vale la pena preguntarse --- y este capítulo tiene motivos para hacerlo --- si un departamento con demanda de lote social documentada tenía tierra fiscal sin destino.

\textbf{La segunda:} los expedientes son de \textbf{2010} y el antecedente de interés turístico es del \textbf{2001}. El proyecto llevaba catorce años en trámite cuando obtuvo la tierra, y ese mismo año obtuvo además la exención impositiva. \textbf{No fue rápido}, y decirlo es parte de describirlo con justicia.
\end{cautela}


\section{Cuando el Estado es el ocupante}

Hay además una vía de adquisición de suelo que no pasa por el sorteo ni por la regularización de ocupantes. La usa el propio Estado, y es la más rápida.

\begin{ficha}{Decreto Nº 1968 del 8 de junio de 2015 --- B.O.\ Nº 19.561}
Expediente 01201.59-198642/2012-0, Ministerio de Educación. \textbf{Declara operada a favor de la Provincia la prescripción adquisitiva del dominio} del inmueble de \textbf{matrícula 3162, Finca Potrero de Gallinato}, departamento La Caldera, de \textbf{3.209,90 m²}, lindante al norte con la Ruta Provincial 11 y en los otros tres rumbos con la matrícula 3161, ``propiedad de María del Valle Esteban y otros''.

\textbf{Destino:} la \textbf{Escuela Nº 4118 ``Dr.\ Gustavo Martínez Zuviría''}, del paraje El Gallinato, emplazada allí \textbf{desde su construcción en el año 1968}.

Fundamento: el Estado provincial detenta la posesión ``en forma pública, pacífica e ininterrumpida, a título de dueño'', y esa posesión ``reconoce origen en la instalación del establecimiento escolar, consentido expresa y pacíficamente en todo tiempo por los sucesivos titulares registrales''.

Encuadre: artículo 4015 del Código Civil y \textbf{Ley nacional 21.477}, modificada por la 24.320. El decreto destaca que se trata de ``un trámite administrativo que \textbf{no requiere intervención judicial}, pues un decreto del Poder Ejecutivo declarando operada la prescripción adquisitiva \textbf{sirve como título suficiente}''.
\end{ficha}

\begin{cautela}
Póngase este procedimiento al lado de los dos que este capítulo describe.

Una familia que ocupa tierra fiscal y quiere el título debe esperar un decreto de adjudicación que, en los casos documentados, tarda entre una y varias décadas; y si va por la vía del sorteo debe ser argentina nativa, acreditar cuatro años de residencia en la localidad y no poseer inmuebles en ningún lugar del mundo.

El Estado provincial, para regularizar su propia ocupación de un terreno privado de tres mil doscientos metros cuadrados, dicta \textbf{un decreto que vale como título} y no pasa por ningún juez.

\textbf{Este libro no objeta el régimen}, que es una ley nacional de 1977 pensada para exactamente esta clase de situaciones: escuelas, puestos sanitarios y comisarías construidos sobre terrenos ajenos con anuencia de los dueños. Es razonable y evita litigios.

Lo que registra es la asimetría del \textbf{procedimiento}: el mismo hecho jurídico --- poseer sin título durante décadas, públicamente y sin oposición --- se resuelve, según quién sea el poseedor, con un decreto de una página o con un expediente de treinta años y requisitos de nacionalidad. La prescripción adquisitiva del ocupante particular exige juicio; la del Estado, no.
\end{cautela}

\begin{cautela}
Dos datos laterales que valen para otros capítulos.

\textbf{El primero:} la escuela está en el paraje \textbf{El Gallinato} desde 1968, sobre la \textbf{Finca Potrero de Gallinato}. El artículo 6º del proyecto de ordenanza del PDUA, analizado en el capítulo \ref{cap:pdua}, enumera entre las unidades del municipio una \textbf{Finca El Gallinato}. Establecer si son la misma unidad, y qué relación guardan las matrículas 3161 y 3162 con aquella enumeración, es parte del antecedente dominial que este libro no obtuvo.

\textbf{El segundo:} el decreto identifica al propietario lindero --- ``María del Valle Esteban y otros'', matrícula 3161 --- y lo hace, otra vez, \textbf{de paso}. Es la tercera vez que este relevamiento averigua quién es dueño de algo en el departamento leyendo un acto que trata de otra cosa.
\end{cautela}


\subsection*{Y el plano apareció}

El apéndice \ref{cap:pedidos} venía pidiendo el \textbf{plano de mensura para prescripción
adquisitiva administrativa Nº 000962}, que es el anexo de este decreto. Ya está, y agrega
cuatro cosas que el texto del decreto no dice.

\begin{ficha}{Plano Nº 000962 --- Dirección General de Inmuebles, mayo de 2014}
\textbf{Mensura para Prescripción Adquisitiva Administrativa.} Departamento La Caldera,
propiedad \textbf{Finca Potrero de Gallinato}. Propietario: \textbf{Esteban María del Valle}.
Pretende prescribir: \textbf{Provincia de Salta}. Matrícula \textbf{3.162}. Título inscripto
en Cédula Parcelaria. Antecedentes gráficos: lámina catastral \textbf{572}.

\textbf{Balance de superficies:} según título \textbf{3.538,08 m\textsuperscript{2}}; a
prescribir \textbf{3.209,90 m\textsuperscript{2}}; remanente \textbf{328,18
m\textsuperscript{2}}. Nota al pie: «la mensura se realiza sobre hechos existentes».

Escala \textbf{1:400}. Registro técnico del 5 de enero de 2015. Firman el Ing.\ Oscar
N.\ Bazón Pouzer, Jefe del Programa Registro Técnico, y el Dr.\ Esteban García Bez, Director.
\end{ficha}

\textbf{Primero: la parcela queda georreferenciada.} El plano trae dos puntos de vinculación:
\textbf{X~3.565.382,30 --- Y~7.271.589,15} y \textbf{X~3.565.448,22 --- Y~7.271.545,27}.

Eso permite una medición que hasta ahora no era posible. El edicto de la Resolución 185/14
publicó las coordenadas de la línea de ribera del río La Caldera sobre la matrícula 4366 ---veintiún pares de los ciento trece del edicto; los otros noventa y dos son del arroyo El Manzano (capítulo \ref{cap:ribera})---, entre \textbf{X~3.561.283 --- Y~7.276.241} y \textbf{X~3.562.295 ---
Y~7.277.843}. Con las dos series en el mismo sistema, la cuenta se hace sola: \textbf{el paraje El Gallinato está a entre seis y siete kilómetros al sudeste} del tramo del río que el capítulo
\ref{cap:ribera} analiza. El croquis de ubicación lo confirma: sitúa
la finca entre el \textbf{río Wierna}, la \textbf{Ruta Nacional 9} y la \textbf{Ruta
Provincial 11}, aguas abajo del puente La Caldera.

\textbf{Segundo: aparece una matrícula más.} El croquis rotula la \textbf{3.160} como
lindera. Con la 3.161 y la 3.162 del decreto, y con la \textbf{3.709} sobre la que el
capítulo \ref{cap:loteo} documenta la aprobación del club de campo Chacras del Gallinato, la Finca Potrero de Gallinato lleva ya dos matrículas identificadas con certeza ---la 3.161 y la 3.162---, una lindera cuya pertenencia a la finca no consta ---la 3.160--- y la 3.709, que podría ser otra, algo que el capítulo \ref{cap:loteo} deja sin establecer. \textbf{Es el mismo
partido de 1909 que siguen los capítulos \ref{cap:fincas} y \ref{cap:siglo}; el segundo lo sigue desde el remate de 1922.}

\textbf{Tercero, y es lo que más importa para este libro: el plano fija un retiro en metros.}
Dos anotaciones lo hacen, y las dos son de organismos viales.

\begin{ficha}{Anotaciones de zona de camino en el Plano 000962}
Recuadro manuscrito de la \textbf{Dirección de Vialidad de Salta}, fechado \textbf{28 de
noviembre de 2014}. Ruta: \textbf{Provincial Nº 11}, principal. Tramo: \textbf{2º El
Desmonte --- Campo Santo}.

Ancho de zona de camino \textbf{declarada o proyectada: 14,70 m}. Ancho de zona de camino
\textbf{según Ley Nº 3.383, art.\ 22 inc.\ a): 60,00 m}.

Sello sobre el propio plano: «\textsc{la d.v.s.\ solicita no ejecutar construcciones de
carácter permanente a una distancia inferior a} \underline{\hspace{1.2cm}}
\textsc{m, medidos desde el eje de la calzada existente}».
\end{ficha}

\begin{cautela}
\textbf{La distancia mínima del sello no se lee con certeza.} Es un número manuscrito de dos
dígitos seguido de «,00», sobre un sello impreso y en una reproducción que no da para más.
Las lecturas posibles son compatibles con la mitad del ancho de zona de camino que el mismo
recuadro declara por ley --- sesenta metros totales, treinta a cada lado del eje ---, pero
\textbf{este libro no lo afirma}: lo que sostiene el argumento que sigue es la existencia de
la cifra y su forma, no su valor exacto. El plano original lo resuelve.
\end{cautela}

\begin{cautela}
\textbf{Y cuarto: lo que esa cifra permite comparar.} Póngase esa cifra al lado del vacío que el capítulo \ref{cap:cot} no pudo llenar.

Para un \textbf{camino provincial}, el aparato administrativo produce tres números exactos
--- el ancho declarado, el ancho legal y la distancia mínima de edificación ---, los
estampa sobre el plano de una parcela de tres mil metros cuadrados, y los hace oponibles a
quien quiera construir.

Para \textbf{ambas márgenes del río}, declaradas no urbanizables por el artículo 43 del
Código de Ordenamiento Territorial, \textbf{no hay ningún número}: hay un trazo en un plano que el municipio publica sólo como imagen y que no resulta mensurable.

\textbf{No es que el Estado no sepa fijar retiros en metros. Los fija cuando quiere.} Es la
comparación más limpia que este relevamiento encontró para lo que el capítulo
\ref{cap:opacidad} describe, y no requiere atribuir intención a nadie: requiere solamente
poner los dos documentos uno al lado del otro.
\end{cautela}

\lamina{lam-plano-000962.jpg}{lam:plano962}%
{Plano de mensura para prescripción adquisitiva administrativa Nº 000962}%
{\textbf{Plano Nº 000962}, Dirección General de Inmuebles, mayo de 2014, escala 1:400. Es la
parcela de 3.209,90 m\textsuperscript{2} de la Escuela Nº 4118 del paraje El Gallinato, que
la Provincia adquirió por el Decreto 1968/15 tras cuarenta y siete años de posesión. Linda al
norte con la Ruta Provincial 11 y en los otros tres rumbos con la matrícula 3.161. Abajo a la
izquierda, el recuadro de zona de camino de la Dirección de Vialidad de Salta, que la lámina
\ref{lam:zonacamino} reproduce en detalle.}%
{Anexo del Decreto Nº 1968 del 8 de junio de 2015, publicado en el Boletín Oficial Nº 19.561
del 16 de junio de 2015. Reproducción facsimilar; la hoja se giró noventa grados para que el
plano quede legible y se recortó el margen en blanco del escaneo.}

\section{Treinta viviendas que necesitan muros}

Antes del Plan Mi Lote hubo obra de vivienda estatal en el departamento, y su documentación dice algo sobre dónde se construyó.

\begin{ficha}{Resolución IPV Nº 170 del 18 de febrero de 2015 --- B.O.\ Nº 19.497}
Expediente 011-0068-271427/2014-0. El Instituto Provincial de Vivienda declara \textbf{apto técnico} y aprueba el pliego para la obra \textbf{``Ejecución de muros de contención para 30 viviendas Vaqueros --- Dpto.\ La Caldera''}.

Presupuesto oficial tope \textbf{\$498.000}, plazo \textbf{sesenta días corridos}, modalidad ajuste alzado, por \textbf{Contratación Directa Nº 25/2014}, con invitación a por lo menos tres empresas del medio.

Firman el interventor Matías Posadas, el jefe del Área de Control de Obras Federico Gauffin y el presidente Marcelo A.\ Ferraris.
\end{ficha}

\begin{cautela}
La obra no es la vivienda: son \textbf{los muros que la sostienen}. Y eso es lo que el documento dice sin decirlo.

Un conjunto de treinta viviendas que requiere una obra específica de contención, contratada aparte y por casi medio millón de pesos, está emplazado en un terreno con pendiente, con riesgo de deslizamiento o con necesidad de nivelación. No es una obra ornamental ni de cierre perimetral: los muros de contención se construyen cuando el suelo, librado a sí mismo, se mueve.

De modo que en el mismo departamento y el mismo año en que este libro documenta clubes de campo con concesiones permanentes de agua y hoteles con exenciones impositivas, el Estado provincial construía \textbf{treinta viviendas sociales en un terreno que necesitaba ser contenido}.

Este libro \textbf{no afirma} que el emplazamiento haya sido inadecuado. Construir con muros de contención en terreno con pendiente es una solución técnica correcta, y en un valle de montaña la pendiente es la regla y no la excepción. Lo que registra es la pregunta que se abre: \textbf{dónde están esas treinta viviendas, sobre qué suelo, y qué pasó con ellas en los temporales de 2018, 2019 y 2026.} Ninguna de las tres respuestas está en este acto.
\end{cautela}


\section{El anexo que faltaba}

Se obtuvo el texto del Anexo de la Resolución 82/16 --- el Reglamento de Asignación de Lotes de Interés Social ---, que este libro venía consignando como no publicado. Estaba publicado: la edición web del Boletín (Nº 19.861, 09/09/2016) lo ofrece como archivo vinculado a la resolución ---una copia certificada de seis folios, escaneada y sin capa de texto, que no aparece en las búsquedas---. Firma el ministro de Infraestructura, Tierra y Vivienda, Baltasar Saravia. Es el marco del que cuelgan las Resoluciones 7/16, 8/16 y 9/16 ya analizadas, y contiene dos artículos que ninguna de ellas anticipaba.

\subsection*{Un registro público que la norma manda llevar}

El artículo 1º enumera las funciones de la autoridad de aplicación. Entre ellas, el inciso e ordena llevar ``un banco de datos actualizado a través del cual se permita evacuar consultas, realizar y responder informes, y efectuar los controles correspondientes''. Y el inciso g es más específico:

\begin{quote}
``Deberá llevar un \textbf{registro público actualizado y de libre acceso} de los Tenedores Precarios y Adjudicatarios en Venta de lotes de interés social.''
\end{quote}

\begin{cautela}
Ese inciso vale por sí solo. Desde 2016 existe una obligación reglamentaria de mantener un registro \textbf{público, actualizado y de libre acceso} de todas las personas que tienen tenencia precaria o adjudicación en venta de lotes de interés social en la provincia.

Este relevamiento no lo encontró. Eso no prueba que no exista --- puede estar en una repartición, en un formato no publicado, o disponible a pedido ---, pero convierte una búsqueda difusa en una pregunta de una línea: \textbf{dónde está el registro que el artículo 1º inciso g manda llevar}, y si es de libre acceso, cómo se accede.

Si existe y es accesible, resuelve de un golpe buena parte de lo que este capítulo persigue lote por lote. Si no existe, es un incumplimiento de una norma que el artículo 11 declara \textbf{de orden público}, no omitible ni modificable bajo pena de nulidad.
\end{cautela}

\subsection*{Y un requisito de nacionalidad}

El artículo 3º fija las condiciones para participar en los sorteos. Su inciso a exige ser mayor de edad y \textbf{``argentino nativo o por opción''}.

\begin{cautela}
Esto es más restrictivo que la Ley 2616 y es de 2016, no de 1951.

La ley de 1951, como este capítulo señaló, coloca a ``los argentinos nativos'' en \textbf{segundo lugar de preferencia}: es un criterio de orden dentro de una lista, no una barrera. El reglamento de 2016, en cambio, convierte la nacionalidad en \textbf{requisito de admisión}: quien no es argentino nativo o por opción no puede siquiera inscribirse al sorteo.

La fórmula ``nativo o por opción'' es además estrecha. Deja fuera a los extranjeros residentes, pero también a los argentinos \textbf{naturalizados} por carta de ciudadanía, que son argentinos a todos los efectos. La tensión con el artículo 20 de la Constitución Nacional --- que reconoce a los extranjeros los mismos derechos civiles que a los ciudadanos, incluido el de poseer bienes raíces --- es evidente y este libro la deja planteada como cuestión jurídica, sin resolverla.
\end{cautela}

Y esto explica por qué ese requisito no alcanza a las adjudicaciones del barrio Juan Manuel de Rosas. La respuesta está en el propio Decreto 63/25 y es estructural, no anómala: esas adjudicaciones \textbf{no salieron de un sorteo}. Se fundaron en el artículo 2º inciso 1 de la Ley 2616, ``por revestir el carácter de actuales ocupantes''. Es la vía de regularización dominial, no la de asignación por concurso, y el requisito de nacionalidad del artículo 3º del reglamento de 2016 rige expresamente ``para poder participar en los sorteos''.

De modo que la provincia tiene dos puertas a la tierra fiscal con reglas de admisión distintas. Por la puerta del sorteo, con puntaje y lista de espera, hay que ser argentino nativo. Por la puerta de la regularización, hay que haber ocupado y no tener otro inmueble ni otra solución habitacional según los registros provinciales. El capítulo ya había mostrado que esas dos vías se ignoran mutuamente en el baremo --- ocupar no suma puntos ---; ahora se ve que también difieren en quién puede entrar.

\subsection*{El resto del reglamento}

\begin{ficha}{Anexo de la Resolución 82/16 --- otros contenidos}
\begin{itemize}[leftmargin=*,nosep]
\item \textbf{Art.\ 3 inc.\ d:} residencia ininterrumpida de \textbf{cuatro años} en la localidad donde se solicita el lote, inmediatamente anterior a la inscripción. Más radicación definitiva en la provincia (inc.\ c) y grupo familiar con al menos un menor a cargo, tutela o curatela (inc.\ b).
\item \textbf{Art.\ 4 --- inhibiciones:} no tener otro trámite en curso; no ser titular registral, cotitular, usufructuario, comprador por boleto ni cesionario de derechos posesorios sobre inmuebles \textbf{en el país o en el exterior}; no haber recibido antes solución habitacional de ningún organismo estatal ni de instituciones privadas.
\item \textbf{Art.\ 5 --- publicidad:} los procedimientos de asignación y sus sorteos son públicos y se publican en la web oficial del Gobierno de manera continua, en el Boletín Oficial y en el diario de mayor difusión ---o el medio más efectivo de la zona--- al menos un día, y en la sede de la Subsecretaría o del municipio; en los cuatro casos, con \textbf{una anticipación mínima de diez días} a la fecha del sorteo. \textbf{Es el único plazo del artículo.}
\item \textbf{Art.\ 6:} los sorteos son públicos y deben estar presididos por escribano público, \textbf{bajo pena de nulidad}.
\item \textbf{Art.\ 7:} el beneficiario es \textbf{tenedor precario} y su obligación principal es habitar el lote de manera inmediata, continua e ininterrumpida. La tenencia garantiza uso y posesión hasta la adjudicación en venta por acto del Poder Ejecutivo.
\item \textbf{Art.\ 8:} prohibición de vender, ceder, transmitir, alquilar o dar en comodato; el infractor pierde \textbf{en forma definitiva} la posibilidad de inscribirse a un nuevo sorteo.
\item \textbf{Arts.\ 9 y 10:} revocación sólo por acto administrativo fundado, previa consulta a las áreas técnicas y legales; impugnación por la Ley 5.348.
\item \textbf{Art.\ 11:} los requisitos y procedimientos son \textbf{de orden público} y no pueden omitirse ni modificarse bajo pena de nulidad, sin perjuicio de que la autoridad exija mayores recaudos.
\end{itemize}
\end{ficha}

El artículo 3 inciso d obliga a matizar --- no a revertir --- la corrección que este capítulo hizo en favor de los funcionarios cuestionados por el acto de Juan Manuel de Rosas.

Se dijo allí que la Resolución 7/16 no contiene una prohibición de adjudicar a habitantes de otros municipios. Eso sigue siendo cierto de esa resolución. Pero el reglamento marco \textbf{sí contiene} una exigencia de cuatro años de residencia ininterrumpida en la localidad donde se solicita el lote, que es una barrera territorial bastante más dura que un certificado policial.

La conclusión, sin embargo, no cambia, y por la misma razón que el requisito de nacionalidad: ese artículo rige \textbf{``para poder participar en los sorteos''}, y las adjudicaciones del Decreto 63/25 no fueron por sorteo sino por ocupación efectiva bajo la Ley 2616. Sigue sin haber, en la normativa examinada, un impedimento legal que alguien haya violado.

Lo que hay es una asimetría que ahora se ve completa. La vía del sorteo exige nacionalidad, cuatro años de arraigo local, no poseer inmuebles en ningún lugar del mundo ---con dos excepciones menores---, no haber recibido nunca ayuda habitacional y no estar incurso en infracciones por ocupación ilegal. La vía de la regularización exige ocupar, y verifica en los registros de Inmuebles y del Instituto Provincial de Vivienda que el ocupante no tenga otro inmueble ni haya recibido otra solución habitacional: la misma inhibición, comprobada de oficio y dentro de la provincia, en lugar de declarada bajo juramento y para todo el mundo. Son dos regímenes de acceso a tierra pública que conviven sin comunicarse, y el más exigente es el que se supone ordenado y transparente.


\section{El Plan Mi Lote}

El Decreto 399/20 crea el Plan destinado a los sectores de mayor vulnerabilidad, con dos fines: entrega en propiedad de lotes y regularización dominial de quienes ocupan sin título. Proyecta \textbf{más de 10.000 lotes} con infraestructura en cuatro años, \textbf{40\% en el área metropolitana y 60\% en el interior}, dotados de agua, luz y apertura de calles. Su implementación se encomienda al Instituto Provincial de Vivienda, con una Unidad Ejecutora Provincial a crear. Su artículo 4º dispone que no podrán ser beneficiarios quienes se encuentren incursos en infracciones vinculadas con la ocupación ilegal de inmuebles públicos o privados.

Cuatro mil lotes proyectados para el aglomerado que integran Salta capital, Vaqueros, La Caldera, San Lorenzo y Cerrillos, en cuatro años. Es la primera magnitud cuantificada del segundo circuito y es contrastable contra la ejecución real.

Dos tensiones, nombradas sin exagerarlas. El artículo 4º excluye a quienes estén incursos en infracciones por ocupación ilegal, y el Decreto 63/25 adjudica a actuales ocupantes: no hay contradicción jurídica --- ocupar no equivale a estar incurso en una infracción, y la regularización de ocupaciones es uno de los dos objetos expresos del plan --- pero la distinción entre uno y otro no está definida en el decreto y la traza la administración caso por caso. Y el plan se asignó al IPV mientras hoy lo ejecuta la Secretaría de Tierras y Bienes del Estado: esa migración debe documentarse.

\section{Estado del aparato y valoración}

De la cadena completa, este libro leyó el texto operativo de casi todo: quedan sin leer los instrumentos marcados como pendientes en la ficha inicial. El anexo de la Resolución 82/16 se obtuvo tarde, y resultó estar en el repositorio del Boletín como imagen escaneada; el de la Ordenanza 437 de Vaqueros, que contiene las exigencias concretas del fraccionamiento, apareció en el digesto en línea del propio municipio (capítulo \ref{cap:vaqueros}).

\begin{cautela}
El caso de la Resolución 82/16 muestra el costo acumulado del patrón que el capítulo \ref{cap:metodo} describió. Es la cabeza de la cadena y a su anexo remiten todas las demás. Y hay un dato que aparece tres veces en la Resolución 9/16, siempre por remisión ---dos de ellas con esta fórmula---: los padrones deben exhibirse ``por un plazo no menor al previsto en el artículo 5º del Anexo de la Resolución Nº 82/16''.

Ese plazo --- cuántos días tiene un vecino para revisar el padrón, detectar un error y presentar una impugnación --- es el número del que depende toda la arquitectura de control ciudadano del sistema. No está en ninguna de las resoluciones publicadas como texto. Leído el anexo ---una copia certificada que el Boletín ofrece sólo como imagen---, el artículo 5º no trata de padrones: regula la publicidad de los sorteos, y su único plazo son \textbf{diez días de anticipación}. La remisión de la Resolución 9/16 sólo puede apuntar a ese número, de modo que \textbf{los padrones deben exhibirse al menos diez días}, pero eso se obtiene por deducción, a partir de una cláusula escrita para otra cosa. El único plazo de exhibición que el régimen escribe con número está en el artículo 7º del Anexo de la Resolución 9/16, y rige para otro padrón: el de asignación directa, que se exhibe cinco días para tachas. No se afirma que sea deliberado. Se afirma que el resultado es verificable y que su forma es precisa: se publican las garantías como texto y el plazo que las hace ejercitables queda en un anexo escaneado.
\end{cautela}

\begin{cautela}
Corresponde la valoración de conjunto, y no es la que este capítulo esperaba escribir.

El régimen de tierra fiscal salteño es, en el papel, considerablemente más riguroso que el régimen municipal de suelo privado que los capítulos \ref{cap:pdua} a \ref{cap:loteo} documentaron. Tiene cupos numéricos explícitos, escribano obligatorio bajo pena de nulidad, padrones publicados, período de tachas, informes socioambientales fundados, baremo escrito, visitas de constatación, derecho de defensa previo a la revocación y recurso administrativo expreso. Nada de eso existe con esa densidad del lado del loteo privado, donde el Código de La Caldera fija un procedimiento que no se aplica y Vaqueros lo tiene fuera de su Código, en las Ordenanzas 437/14 y 535/18, sin que conste tampoco su aplicación.

Al solicitante de tierra fiscal se le exige certificado policial de residencia, declaración jurada bajo pena de falso testimonio, referencias vecinales, franja horaria de visita domiciliaria y construcción en tres meses bajo apercibimiento. Al desarrollador que lotea sobre una microcuenca condenada judicialmente se le exige un certificado de aptitud ambiental municipal cuya existencia este libro no pudo verificar para el loteo que el amparo litiga; el único certificado ambiental que encontró en el corredor es provincial, y lo presentó un loteador al ser intimado por el municipio (capítulo \ref{cap:amparo}).

No hay en la tierra fiscal salteña un régimen laxo que favorezca la discrecionalidad: hay un régimen exigente aplicado a quien menos tiene y un régimen tenue aplicado a quien más puede. Los dos circuitos de suelo del departamento La Caldera no se distinguen por su legalidad sino por la intensidad con que el Estado se ocupa de cada uno.
\end{cautela}
